\documentclass[11pt, a4paper]{article}
\usepackage[a4paper, margin=2cm]{geometry}
\usepackage{fontspec}
\usepackage[english, russian, bidi=basic, provide=*]{babel}
\babelprovide[import, onchar=ids fonts]{russian}
\babelfont{rm}{Noto Sans}
\usepackage{amsmath, amssymb, amsthm}
\usepackage{booktabs}
\usepackage{enumitem}
\usepackage{hyperref}

\title{\textbf{Полные решения заданий}\\ \large Задание 4. Линейные операторы, квадратичные формы и евклидовы пространства}
\author{Линейная алгебра}
\date{2024}

\begin{document}

\maketitle

% ==========================================
\section{Раздел 1. Линейные операторы и их матрицы}
% ==========================================

\subsection*{Задание 1.1}
\textbf{Условие:} Оператор $A$ поворачивает каждый вектор плоскости $\mathbb{R}^2$ на угол $\phi = \frac{\pi}{3}$ (60$^\circ$) против часовой стрелки. Найдите матрицу оператора $A$ в стандартном базисе.

\textbf{Решение:}
Матрица поворота на угол $\phi$ против часовой стрелки в каноническом базисе имеет вид:
$$ M_A = \begin{pmatrix} \cos\phi & -\sin\phi \\ \sin\phi & \cos\phi \end{pmatrix} $$
Подставим значение $\phi = \frac{\pi}{3}$:
$$ \cos\left(\frac{\pi}{3}\right) = \frac{1}{2}, \quad \sin\left(\frac{\pi}{3}\right) = \frac{\sqrt{3}}{2} $$
Следовательно, матрица оператора равна:
$$ M_A = \begin{pmatrix} \frac{1}{2} & -\frac{\sqrt{3}}{2} \\ \frac{\sqrt{3}}{2} & \frac{1}{2} \end{pmatrix} $$

\textbf{Ответ:} $M_A = \begin{pmatrix} 0.5 & -\frac{\sqrt{3}}{2} \\ \frac{\sqrt{3}}{2} & 0.5 \end{pmatrix}$.

---

\subsection*{Задание 1.2}
\textbf{Условие:} Линейный оператор $A$ переводит базисные векторы $e_1, e_2$ в $A(e_1) = (2, 5)$, $A(e_2) = (-1, 3)$. Найдите образ вектора $x = (3, -2)^T$ под действием оператора $A$.

\textbf{Решение:}
1. Составим матрицу оператора $A$, расположив образы базисных векторов по столбцам:
$$ A = \begin{pmatrix} 2 & -1 \\ 5 & 3 \end{pmatrix} $$
2. Вычислим результат применения оператора к вектору $x$:
$$ A(x) = A \cdot x = \begin{pmatrix} 2 & -1 \\ 5 & 3 \end{pmatrix} \begin{pmatrix} 3 \\ -2 \end{pmatrix} = \begin{pmatrix} 2\cdot 3 + (-1)\cdot(-2) \\ 5\cdot 3 + 3\cdot(-2) \end{pmatrix} = \begin{pmatrix} 6 + 2 \\ 15 - 6 \end{pmatrix} = \begin{pmatrix} 8 \\ 9 \end{pmatrix} $$

\textbf{Ответ:} $A(x) = (8, 9)^T$.

\newpage

% ==========================================
\section{Раздел 2. Квадратичные формы}
% ==========================================

\subsection*{Задание 2.1}
\textbf{Условие:} Составьте матрицу квадратичной формы $Q(x_1, x_2, x_3) = 3x_1^2 + 5x_2^2 - 2x_3^2 + 4x_1x_2 - 6x_1x_3 + 2x_2x_3$.

\textbf{Решение:}
Элементы симметричной матрицы квадратичной формы $A = (a_{ij})$ формируются следующим образом:
\begin{itemize}
\item На главной диагонали $a_{ii}$ стоят коэффициенты при $x_i^2$.
\item Внеглавные элементы $a_{ij} = a_{ji}$ равны половине коэффициента при произведении $x_i x_j$.
\end{itemize}

Вычисляем элементы:
\begin{align*}
a_{11} &= 3, \quad a_{22} = 5, \quad a_{33} = -2, \\
a_{12} &= a_{21} = \frac{4}{2} = 2, \\
a_{13} &= a_{31} = \frac{-6}{2} = -3, \\
a_{23} &= a_{32} = \frac{2}{2} = 1.
\end{align*}

\textbf{Ответ:} $A = \begin{pmatrix} 3 & 2 & -3 \\ 2 & 5 & 1 \\ -3 & 1 & -2 \end{pmatrix}$.

---

\subsection*{Задание 2.2}
\textbf{Условие:} Исследуйте на знакоопределенность квадратичную форму $Q(x_1, x_2) = 2x_1^2 + 6x_1x_2 + 5x_2^2$ с помощью критерия Сильвестра.

\textbf{Решение:}
1. Составим матрицу квадратичной формы:
$$ A = \begin{pmatrix} 2 & 3 \\ 3 & 5 \end{pmatrix} $$
2. Вычислим угловые миноры:
\begin{itemize}
\item $\Delta_1 = |2| = 2 > 0$
\item $\Delta_2 = \begin{vmatrix} 2 & 3 \\ 3 & 5 \end{vmatrix} = 2 \cdot 5 - 3 \cdot 3 = 10 - 9 = 1 > 0$
\end{itemize}
Так как все угловые главные миноры строго положительны ($\Delta_1 > 0, \Delta_2 > 0$), по критерию Сильвестра квадратичная форма является **положительно определенной**.

\textbf{Ответ:} Квадратичная форма положительно определена.

\newpage

% ==========================================
\section{Раздел 3. Ортогонализация Грама-Шмидта}
% ==========================================

\subsection*{Задание 3.1}
\textbf{Условие:} Примените процесс ортогонализации Грама-Шмидта к системе векторов $f_1 = (1, 1, 0)$ и $f_2 = (1, 0, 1)$.

\textbf{Решение:}
1. Первым ортогональным вектором берем $u_1 = f_1 = (1, 1, 0)$.
2. Вектор $u_2$ ищем в виде $u_2 = f_2 - \frac{(f_2, u_1)}{\|u_1\|^2} u_1$:
   \begin{itemize}
   \item Скалярное произведение: $(f_2, u_1) = 1\cdot 1 + 0\cdot 1 + 1\cdot 0 = 1$.
   \item Квадрат нормы: $\|u_1\|^2 = 1^2 + 1^2 + 0^2 = 2$.
   \item Коэффициент проекции: $\frac{(f_2, u_1)}{\|u_1\|^2} = \frac{1}{2}$.
   \end{itemize}
Вычисляем $u_2$:
$$ u_2 = (1, 0, 1) - \frac{1}{2}(1, 1, 0) = \left(1 - \frac{1}{2}, \; 0 - \frac{1}{2}, \; 1 - 0\right) = \left(\frac{1}{2}, \; -\frac{1}{2}, \; 1\right). $$

\textbf{Проверка ортогональности:}
$$ (u_1, u_2) = 1 \cdot \frac{1}{2} + 1 \cdot \left(-\frac{1}{2}\right) + 0 \cdot 1 = \frac{1}{2} - \frac{1}{2} + 0 = 0. $$

\textbf{Ответ:} Ортогональные векторы: $u_1 = (1, 1, 0)$, $u_2 = (0.5, -0.5, 1)$.

\newpage

% ==========================================
\section{Раздел 4. Бонусные задания и Теория}
% ==========================================

\subsection*{Бонус 1}
\textbf{Условие:} Докажите, что собственные векторы симметричной матрицы, отвечающие различным собственным значениям, ортогональны.

\textbf{Доказательство:}
Пусть $A = A^T$ --- вещественная симметричная матрица, и пусть $v_1, v_2$ --- ее собственные векторы с различными собственными значениями $\lambda_1 \ne \lambda_2$:
$$ A v_1 = \lambda_1 v_1, \quad A v_2 = \lambda_2 v_2. $$
Рассмотрим скалярное произведение $(A v_1, v_2)$:
$$ (A v_1)^T v_2 = (v_1^T A^T) v_2 = v_1^T (A v_2) = v_1^T (\lambda_2 v_2) = \lambda_2 (v_1^T v_2) = \lambda_2 (v_1, v_2). $$
С другой стороны, подставляя $A v_1 = \lambda_1 v_1$:
$$ (A v_1)^T v_2 = (\lambda_1 v_1)^T v_2 = \lambda_1 (v_1^T v_2) = \lambda_1 (v_1, v_2). $$
Приравниваем оба выражения:
$$ \lambda_1 (v_1, v_2) = \lambda_2 (v_1, v_2) \implies (\lambda_1 - \lambda_2) (v_1, v_2) = 0. $$
Так как по условию $\lambda_1 \ne \lambda_2$, разность $(\lambda_1 - \lambda_2) \ne 0$. Следовательно:
$$ (v_1, v_2) = 0, $$
что и означает ортогональность векторов $v_1$ и $v_2$. \qed

---

\subsection*{Теория}

\textbf{a) Сформулируйте критерий Сильвестра для отрицательно определенной квадратичной формы.}
\begin{proof}[Ответ]
Квадратичная форма $Q(x)$ с матрицей $A$ является отрицательно определенной тогда и только тогда, когда знаки ее главных угловых миноров чередуются, начиная с минуса:
$$ (-1)^k \Delta_k > 0 \quad \text{для всех } k = 1, 2, \dots, n, $$
то есть: $\Delta_1 < 0, \; \Delta_2 > 0, \; \Delta_3 < 0, \; \Delta_4 > 0$ и так далее.
\end{proof}

\textbf{b) Какой оператор называется самосопряженным (симметричным)?}
\begin{proof}[Ответ]
Оператор $A$, действующий в евклидовом пространстве, называется самосопряженным, если для любых векторов $x, y$ выполняется равенство скалярных произведений:
$$ (A x, y) = (x, A y). $$
В ортонормированном базисе матрица самосопряженного оператора является симметричной ($A = A^T$).
\end{proof}

\end{document}